% =====================================================================
%  24.5  The Engine Boundary: when centroid decomposition is the WRONG tool.
%  WORKED EXAMPLES ONLY -- NO \solvedproblem (in-section count stays 4).
%  Exports: sec:c18-boundary  def:c18-map  lst:c18-map
%  DISCHARGES (prose + \chref/\secref only, NO re-link):
%    * Ch.22 (c16) \secref{c16-boundary}: "path COUNTING / distance stats is a
%      centroid job, not HLD" -> answered by \secref{c18-count}/\secref{c18-index}.
%    * Ch.23 (c29) "over ALL paths -> centroid" forward-cite -> \secref{c18-beyond}.
%    * Ch.16 (c11) LCA: the pairwise oracle the centroid index generalises.
%  Stays OUT (name + forward-cite, no code): single fixed-path RANGE updates ->
%    HLD Ch.22 (c16); subtree multiset unions -> DSU-on-tree Ch.17 (c28); a state
%    per node -> tree DP / rerooting Ch.13 (c10) / Ch.20 (c20); offline reorderable
%    node-set queries -> Mo Ch.27 (c21); dynamic topology -> link-cut Ch.25 (c23).
%  Do NOT \term anything already indexed ("centroid", "centroid tree" -- Ch.12 /
%    this chapter's 00/01).  This section uses ZERO \term.
% =====================================================================
\section{The Engine Boundary: When Centroid Is the Wrong Tool}\label{sec:c18-boundary}

\begin{previouslybox}[Four sections built one engine; this one says when to leave it]
Sections 24.1--24.4 made the centroid decomposition do a great deal: build the
\(\bigO{\log n}\)-deep centroid tree, count pairs by distance, index a component for
nearest-marked queries, and fold a product over colour-balanced paths. It is seductive
for exactly that reason --- once you can ``split at the balance and recurse,'' every
tree problem with a path in it starts to look centroid-shaped. But Chapter~22 (\chref{c16})
owns path \emph{ranges}, Chapter~17 (\chref{c28}) owns subtree-set \emph{unions},
Chapter~13 (\chref{c10}) owns per-node \emph{DP states}, and Chapter~25 (\chref{c23})
owns \emph{changing} trees --- and each of those beats centroid decomposition at its own
job, often by a full \(\log n\). A boundary section earns its place because the wrong
pick here is rarely a wrong \emph{answer}; it is a \(\tle\) or a needlessly \(n\log^2 n\)
where an \(n\log n\) or even \(n\) was available.
\end{previouslybox}

\begin{topicintro}
Three questions put any tree problem on the map below. \emph{What is the unit of work?}
A single rooted subtree is a DP (\(\chref{c10}\)); a chosen path is a range (HLD,
\(\chref{c16}\)); a whole set of marked nodes is an index (centroid, \(\secref{c18-index}\));
\emph{all} simple paths is centroid territory (\(\secref{c18-count}\), \(\secref{c18-beyond}\)).
\emph{Is the tree fixed?} Centroid decomposition is built once; if edges appear and
vanish between queries, it is invalid --- that is link-cut (\(\chref{c23}\)). \emph{Do you
merge child sets or split the component?} Bottom-up multiset merging is DSU-on-tree /
sack (\(\chref{c28}\)); top-down splitting by balance is centroid. The tell is directional:
sack \emph{accumulates} children into a parent, centroid \emph{dissolves} a component at
its middle.

\textbf{Prerequisites.} The whole of 24.1--24.4; Chapter~22 (\(\secref{c16-boundary}\)) and
Chapter~23 (\(\secref{c29-boundary}\)) for the two neighbouring boundary maps that
forward-cited here; Chapter~16 (\chref{c11}) for the pairwise \(\dist\) oracle.
\end{topicintro}

% ---------------------------------------------------------------
\subsection{The axis that decides the engine: what the query ranges over}
% ---------------------------------------------------------------

\begin{definition}{The tree-path engine map}{c18-map}
Classify by the \emph{set the query quantifies over} and whether the \emph{topology is
fixed}. \(\{\text{all simple paths},\ \text{property additive along path},\ \text{fixed tree}\}\)
is centroid decomposition (this chapter). \(\{\text{one chosen path},\ \text{range
update/query},\ \text{fixed tree}\}\) is heavy-light (\(\chref{c16}\)).
\(\{\text{subtree},\ \text{merge child multisets}\}\) is DSU-on-tree (\(\chref{c28}\)).
\(\{\text{per node},\ \text{local state recurrence}\}\) is tree DP / rerooting
(\(\chref{c10}\), \(\chref{c20}\)). \(\{\text{arbitrary node sets},\ \text{offline,
reorderable}\}\) is Mo on trees (\(\chref{c21}\)). \(\{\text{path or subtree},\ \text{edges
change}\}\) is link-cut (\(\chref{c23}\)). The cells below name the tell that selects one.
\end{definition}

\begin{center}
\small
\renewcommand{\arraystretch}{1.26}
\setlength{\tabcolsep}{4pt}
\begin{tabular}{@{}p{0.30\linewidth} p{0.24\linewidth} >{\raggedright\arraybackslash}p{0.24\linewidth} >{\centering\arraybackslash}p{0.14\linewidth}@{}}
\toprule
\textbf{Symptom in the spec} & \textbf{Right engine} & \textbf{Why centroid loses} & \textbf{Where}\\
\midrule
count/aggregate over \emph{all} paths, value additive along the path & centroid decomposition & --- it wins; this is its home & \(\secref{c18-count}\)/\(\secref{c18-beyond}\)\\
nearest / farthest marked node, online set, fixed tree & centroid distance index & --- it wins; \(\bigO{\log n}\) per touch & \(\secref{c18-index}\)\\
update and query a \emph{single} path, op is a monoid & heavy-light + segtree & centroid reprocesses every node per query, \(\bigO{n}\); no per-path laziness & Ch.~22 \(\chref{c16}\)\\
answer for each \emph{subtree} a set statistic (distinct colours, \(k\)-th) & DSU-on-tree / sack & centroid splits by balance, not by subtree; sack keeps \(\bigO{n\log n}\) merge & Ch.~17 \(\chref{c28}\)\\
a \emph{state per node} from a local recurrence (independent set, matching) & tree DP / rerooting & no distance / all-pairs dimension to exploit; a fold is enough & Ch.~13 \(\chref{c10}\)\\
many \emph{offline} queries over node sets, order free, no inverse & Mo on trees & centroid is online but has no handle on an arbitrary node subset & Ch.~27 \(\chref{c21}\)\\
path/subtree queries but edges are \emph{linked and cut} between ops & link-cut trees & the centroid tree is built once and dies when topology moves & Ch.~25 \(\chref{c23}\)\\
over all paths, but the per-node value is a line envelope & centroid bucket \emph{hosting} Li Chao & plain centroid folds scalars; combine the two & Ch.~23 \(\chref{c29}\)\\
\bottomrule
\end{tabular}
\end{center}

\begin{lstlisting}[language={}, basicstyle=\ttfamily\small, numbers=none, frame=single,
  caption={Decision ladder for tree-path problems. Read top to bottom and stop at the first branch that matches. Hands off to Chapters 13/17/22/25/27 whenever a different engine already owns the cell; centroid decomposition is the LAST resort for "all paths," not the first for "a path."}, label={lst:c18-map}]
  What does the query RANGE over?

  +-- A SINGLE chosen path, and you UPDATE + QUERY it (sum/min/assign on a subpath):
  |     -> HEAVY-LIGHT DECOMPOSITION + a segment tree (Ch.22, c16).
  |        Centroid is wrong: it has no notion of "the nodes on path u--v only"
  |        and would touch O(n) buckets to answer one path.

  +-- A SUBTREE (everything under some root), usually a SET statistic:
  |     -> DSU-ON-TREE / SACK (Ch.17, c28) or Euler-tour + Fenwick (Ch.18, c14).
  |        Centroid splits the tree at its CENTRE, top-down; sack merges children
  |        bottom-up.  "merge the children's sets" is sack's signature, not centroid's.

  +-- A STATE PER NODE from a LOCAL recurrence (max independent set, tree matching):
  |     -> TREE DP (Ch.13, c10); changing the root -> REROOTING (Ch.20, c20).
  |        There is no all-pairs / distance term, so halving buys nothing -- a
  |        single O(n) DFS fold already solves it.

  +-- A WHOLE SET of marked nodes and you ask DISTANCE-TO-SET, online:
  |     -> CENTROID DISTANCE INDEX (this ch., 24.3).  Mark/query walk the O(log n)
  |        centroid ancestors of one node.  (Fixed tree only -- see the OUT box.)

  +-- ALL simple paths, with a value that ADDS along the path and a filter that is
  |   a threshold / congruence / dominance on that sum:
  |     -> CENTROID PAIR COUNTING (24.2) or the general D&C-on-trees template with
  |        an ordered structure in the bucket (24.4).  Claim each path at its top
  |        centroid; a rooted DP cannot see paths that cross its root.

  +-- MANY OFFLINE queries over arbitrary node sets, you may REORDER them, aggregate
  |   has no inverse:
  |     -> MO ON TREES (Ch.27, c21).  Centroid answers online but cannot index an
  |        unconstrained subset of nodes chosen by the query.

  STILL OUT OF SCOPE for centroid decomposition (forward-cites only, no code):
    * edges LINKED / CUT between operations -> LINK-CUT TREES (Ch.25, c23); the
      centroid tree is computed once on a FIXED topology and is invalid after an edit.
    * the aggregate over a path is an ENVELOPE of lines queried at a variable x ->
      LI CHAO on a tree (Ch.23, c29); you may still host it in a centroid bucket, but
      the engine doing the work is Li Chao, not the balance split.
\end{lstlisting}

% ---------------------------------------------------------------
\subsection{Worked boundary calls (read a spec, name the engine)}
% ---------------------------------------------------------------

\begin{comparisonbox}[Five specifications, five different answers -- none solved here]
\textbf{(A) ``\(n\le10^5\), \(K\le n\); count unordered pairs at distance exactly \(K\).''}
The unit of work is \emph{all} pairs and the value (distance) is additive along the path.
Chapter~13's \(\bigO{nK}\) subtree DP times out because \(K\) is not small. This is
\(\secref{c18-count}\): claim each pair at its top centroid and read a depth histogram ---
\(\bigO{n\log n}\), independent of \(K\). Centroid is the right call.

\textbf{(B) ``fixed tree, node values; `add \(v\) to every node on path \(u\)--\(w\)', `query the sum on path \(p\)--\(q\)','' }
A single chosen path is \emph{updated and queried}, and the operation (add, sum) is a
monoid that composes into a lazy tag. Heavy-light plus a lazy segment tree
(Ch.~22 \(\chref{c16}\), over Ch.~20 \(\chref{c20}\)) is the answer. Centroid is
\emph{wrong}: it cannot restrict its work to the nodes of one path, so it degrades to
touching \(\bigO{n}\) buckets per query with no laziness to amortise them.

\textbf{(C) ``\(n\le10^5\) coloured tree; \(q\) offline queries `how many distinct colours in the subtree of \(v\)'.''}
The domain is a \emph{subtree} and the statistic merges the children's colour sets
bottom-up --- sack's signature (Ch.~17 \(\chref{c28}\)), or an Euler-tour interval plus a
Fenwick (Ch.~18 \(\chref{c14}\)). Centroid splits at the centre, not at a subtree
boundary, so it has no natural handle on "the descendants of \(v\)." Wrong tool.

\textbf{(D) ``operations `paint node black' and `distance from \(v\) to the nearest black node', online, fixed tree.''}
This is the \emph{nearest-marked} pattern (the classic \emph{Xenia and Tree}; its
statement is Ch.~12's, \chref{c09}). A per-centroid distance table (\(\secref{c18-index}\))
answers it in \(\bigO{\log n}\) per operation --- centroid is exactly right. Note the
contrast with (B): both are "online on a fixed tree," yet (B) is a path \emph{range} (HLD)
while (D) is a distance-to-\emph{set} (centroid). The unit of work, not the online-ness,
picks the engine.

\textbf{(E) ``\(n\le10^5\); `link \(u\)--\(v\)', `cut the edge at \(v\)', `query path aggregate', interleaved.''}
The topology moves. Any decomposition built once --- centroid \emph{or} heavy-light --- is
invalid after an edit. This is link-cut trees (Ch.~25 \(\chref{c23}\)). It is the sharpest
boundary this section draws: centroid decomposition assumes a \emph{fixed} tree, full stop.
\end{comparisonbox}

\begin{observation}[The all-pairs tell, and why a fixed root blinds a DP]
A rooted tree DP at \(r\) can fold anything \emph{inside} a subtree, but a path whose two
endpoints lie in \emph{different} child-subtrees of \(r\) passes through \(r\) and is never
seen by a single bottom-up pass --- and if it does not pass through \(r\), it lives in one
child and the DP misses it entirely unless it recurses anyway. Centroid decomposition
fixes the view instead of rooting it: the top-centroid claim (\(\defref{c18-claim}\)) makes
every pair appear exactly once, at the centroid that separates it. That is the whole reason
"quantifies over all paths" maps here and "quantifies over one subtree" maps to Ch.~13.
\end{observation}

\begin{notebox}[Centroid versus heavy-light: the honest tie-break]
Both \(\bigO{\log n}\)-decompose a tree, and beginners conflate them. They answer
\emph{orthogonal} questions. Heavy-light cuts a tree into root-to-leaf \emph{chains} so
that any single path \(u\)--\(v\) becomes \(O(\log n)\) contiguous array ranges --- it is a
\emph{path-as-range} engine, perfect for update-and-query on one path with a lazy
structure (Ch.~22). Centroid cuts at the \emph{balance point} so that any node becomes
\(O(\log n)\) ancestors across \emph{whole components} --- a \emph{component-as-index}
engine, perfect for all-pairs counts and distance-to-a-set. If the spec says ``the path
\(u\)--\(v\),'',' think heavy-light; if it says ``all pairs'' or ``nearest such node,'','
think centroid. Neither replaces the other, and (B) versus (D) above is the exercise that
proves it.
\end{notebox}

\begin{stuckbox}[If you reach for centroid but the tree has \(n\le 2000\)]
Small \(n\) is a licence to brute force. An all-pairs problem on \(n\le2000\) is a
\(\bigO{n^2}\) double loop with \(\dist\) from Ch.~16's \(\bigO{1}\) LCA oracle --- no
decomposition, no histogram, far less to get wrong. Reach for the centroid engine only when
the \(\bigO{n^2}\) or \(\bigO{nK}\) rival genuinely does not finish; a \(\bigO{n\log^2 n}\)
template that you mis-remember is worse than the correct quadratic you can write in three
lines. The cost of the technique is real code (deleted markers, per-arm subtraction, a
bucket structure), and it is only amortised by scale.
\end{stuckbox}

\begin{recallbox}[24.5 in one breath]
\emph{Pick by what the query ranges over, not by how ``tree-y'' it feels:} one chosen path
with updates \(\to\) heavy-light (Ch.~22); a subtree's set \(\to\) DSU-on-tree (Ch.~17) or
Euler-tour (Ch.~18); a per-node local state \(\to\) tree DP (Ch.~13); arbitrary offline node
sets \(\to\) Mo (Ch.~27); moving edges \(\to\) link-cut (Ch.~25). Centroid owns the two cells
nobody else can: \emph{all} paths with an additive value (24.2/24.4) and distance-to-a-set on
a fixed tree (24.3). Discharges the forward-cites from Ch.~22 (\(\secref{c16-boundary}\))
and Ch.~23 (\(\secref{c29-boundary}\)) that punted "all-pairs / counting paths" to this
chapter --- those promises now sit in \(\secref{c18-count}\), \(\secref{c18-index}\) and
\(\secref{c18-beyond}\).
\selftest{\item A query asks the maximum on path \(u\)--\(v\) after point updates.  Centroid
or heavy-light? \emph{(Heavy-light: a single path is being queried, and the max is a composable
monoid; centroid has no per-path laziness.)}
\item Same tree, but the query is ``distance from \(v\) to the closest marked node,'' marks
change online.  Why does the answer switch? \emph{(Now the set of marked nodes is the object,
not a path --- the centroid distance index of 24.3 answers in \(\bigO{\log n}\) per touch.)}
\item Edges are added and removed between queries.  Why is centroid decomposition disqualified
entirely? \emph{(It is built once on a fixed topology; any link/cut invalidates every
\(\mathit{par}_c\) table --- that is link-cut territory, Ch.~25.)}}
\end{recallbox}
